\documentclass[11pt]{article}

\setlength{\parindent}{0pt}
\usepackage{xltxtra}
\usepackage{hyperref}
\hypersetup{hidelinks}
\usepackage{url}
\urlstyle{tt}
\usepackage{xcolor}
\definecolor{CVBlue}{RGB}{23,110,191}
\usepackage{calc}
\usepackage{graphicx}
\usepackage{tikz}
\usetikzlibrary{calc}
\usepackage{fontspec}
\usepackage{xeCJK}
\usepackage{makecell}
\usepackage{enumitem}

\setmainfont[
    Path = fonts/Main/,
    Extension = .otf,
    BoldFont = texgyretermes-bold.otf,
]{texgyretermes-regular.otf}

\setCJKmainfont[
    Path = fonts/hansans/,
    Extension = .ttf,
    BoldFont = NotoSansSC-Bold.ttf,
]{NotoSansSC-Regular.otf}

\usepackage{relsize}
\usepackage{xspace}
\usepackage{fontawesome}

\usepackage[
    a4paper,
    left=1.0cm,
    right=1.0cm,
    top=1.2cm,
    bottom=0.8cm,
    nohead
]{geometry}

\linespread{1.15}

\usepackage{titlesec}
\setlist{noitemsep, parsep=0pt}
\setlist[itemize]{topsep=0.1em, leftmargin=*, itemsep=0pt}

\newlength{\interProjectSpacing}
\setlength{\interProjectSpacing}{0.4em}
\newcommand{\projectsep}{\vspace{\interProjectSpacing}}

\newlength{\intraProjectTitleSep}
\setlength{\intraProjectTitleSep}{0.15em}
\newcommand{\titlebreak}{\\[\intraProjectTitleSep]}

\titleformat{\section}
{\large\bfseries\raggedright}
{}{0em}
{}
[{\color{CVBlue}\titlerule}]
\titlespacing*{\section}{0cm}{0.5em}{0.3em}

\begin{document}
\pagenumbering{gobble}


\begin{tikzpicture}[remember picture, overlay]
    \node[anchor = north west] at ($(current page.north west)+(0.3cm,+0.8cm)$) {\includegraphics[height=5cm]{zju.png}};
\end{tikzpicture}%

\centerline{\LARGE\bfseries{曹天赐}}
\centerline{\normalsize{\faPhone\ 15556927015 \quad \faWeixin\ 15556927015 \quad \faEnvelopeO\ \href{ziyectc@gmail.com}{ziyectc@gmail.com} \quad \faGithub\ \href{https://github.com/ziyefbk}{ziyefbk}}}

\section{教育背景}
浙江大学 控制科学与技术学院 \hfill 2023.9 -- 至今\\
\textbf{竺可桢学院 工程教育高级班} · 自动化（控制）专业 \hfill 大三\\
GPA: 3.93/4.0

\section{科研与项目经历}

\textbf{实验室科研} \hfill 2024.06 -- 至今 \titlebreak
参与CCF-A类编程语言顶级会议论文（在投）\\
\textit{Formalizing the Linux eBPF Core ISA: A Mechanized Operational Semantics and Its Real-World Applications}\\
研究方向：形式化方法，程序语言理论，ebpf，Rocq

\section{项目经历}

\textbf{Linux内核贡献} \hfill 2025.01 -- 至今 \titlebreak
为Linux内核提交1个patch（\href{https://github.com/kernel-patches/bpf/commit/9d21199842247ab05c675fb9b6c6ca393a5c0024}{commit 9d21199}），提高BPF verifier数值追踪精度\\
共同参与patch提交 div \& mod analysis（\href{https://github.com/kernel-patches/bpf/commits/bpf-next_base}{bpf-next\_base}）\projectsep

\textbf{Tnum与Wint数值抽象的Rust实现} \hfill 2025.07 -- 2025.10 \titlebreak
开源项目复现：ISSTA'25 \& TOPLAS'15，将数值抽象域C++实现迁移至Rust，设计Rust crate支持静态分析与程序验证（\href{https://github.com/ziyefbk/solana-sbpf/tree/main}{项目链接}）

\projectsep

\textbf{2025 iGEM 国际遗传工程机器大赛} \hfill 2025.01 -- 2025.11 \titlebreak
ZJU-China（浙江大学代表队）· 银奖（Silver Medal），负责硬件设计与开发（\href{https://2025.igem.wiki/zju-china/hardware}{硬件设计}）：传感器选型、电路设计、硬件3D建模与硬件调试，同时负责部分网页设计；与生物学、计算机科学团队跨学科协作（\href{https://2025.igem.wiki/zju-china/}{项目主页}）

\projectsep

\textbf{基于生物电路的重金属检测系统} \hfill 2025.01 -- 2025.11 \titlebreak
2025 iGEM国际遗传工程机器大赛参赛项目（ZJU-China浙江大学代表队），获银奖（Silver Medal）\\
担任硬件设计与开发负责人：传感器选型、电路设计、硬件3D建模与调试\\
主导跨学科协作，与生物学、计算机科学团队共同完成项目（\href{https://2025.igem.wiki/zju-china/hardware}{硬件设计}）

\projectsep

\textbf{嵌入式AI系统与模型部署工程实践} \hfill 2024.06 -- 至今 \titlebreak
融合嵌入式开发与深度学习，完成多个端到端AI系统项目\\
\textbf{魔杖动作识别控制系统}：基于ESP32与MPU6050构建传感器数据采集系统，训练CNN/Transformer混合模型完成姿态识别，驱动电机控制机器狗完成指定动作指令\\
\textbf{嵌入式AI系统}：自学Stanford CS231n与MIT TinyML，完成CNN、ResNet、RNN/LSTM、Transformer等模型实现；掌握模型量化、剪枝、知识蒸馏等优化技术，部署至微控制器平台

\section{专业技能}
\begin{itemize}[nosep, itemsep=0.1em]
    \item \textbf{编程语言：} C/C++, Python, Rust, MATLAB
    \item \textbf{AI/机器学习框架：} PyTorch
    \item \textbf{嵌入式开发：} STM32, ESP32, Arduino, 51, 电路设计与PCB
    \item \textbf{开发工具：} \LaTeX, Simulink, Git, Linux
    \item \textbf{专业基础：} 微积分、线性代数、矩阵论、概率论与数理统计、数据结构、机器学习与深度学习、信号分析与处理、自动控制原理
\end{itemize}

\section{获奖情况}
\begin{itemize}[itemsep=0.2em]
    \item \textbf{2025 iGEM 国际遗传工程机器大赛} 银奖
    \item \textbf{浙江大学三等奖学金} \hfill 2023-2024学年
    \item \textbf{浙江大学二等奖学金} \hfill 2024-2025学年
\end{itemize}

\section{其他}
\begin{itemize}[itemsep=0.1em]
    \item \textbf{研究兴趣：} 机器学习、MLsys、形式化方法、程序语言理论
    \item \textbf{英语水平：} CET-6 512分
\end{itemize}
\end{document}
